% Generated by scripts/txt2tex.py from chapters/chapter_02.txt. Edit the TXT, not this file.
\chapterlayout{02}
\chapterquestion{איך הופכות אלפי דוגמאות לידע, ומה קורה כשהן מטעות?}
\chapter{מנתונים ללמידה}\label{ch:02}
\section{להריץ את הסרט אחורה}
השאלה על הבצק עדיין מהבהבת על המסך. אבל כדי להבין מאיפה תגיע התשובה, צריך להריץ את הסרט אחורה, הרבה לפני הערב הזה. עוד לפני שמישהו ידע שהקוראת תשאל, המערכת כבר עברה מסע ארוך של התנסות בדוגמאות.
גם הקוראת לא נולדה אופה. כשהיא נזכרת איך למדה, היא לא רואה מול עיניה ספר הוראות, אלא שורה ארוכה של ניסיונות: כיכר שיצאה שטוחה, עוגה שהתפוררה, ופעם אחת, סוף סוף, לחם שכולם ביקשו ממנו עוד. אף ניסיון לא לימד אותה כלל מפורש. ביחד הם בנו אצלה תחושה. ומה שהיא לא שמה לב אליו אז: גם הניסיונות עצמם נבחרו בשבילה. היא אפתה רק במטבח שלה, בתנור שלה, בקמח מהמכולת הקרובה.
מערכת לומדת עוברת מסלול דומה, אבל כל צעד בו הוא החלטה של בני אדם. את המערכת הגדולה שהקוראת משוחחת איתה קשה לפתוח, ולכן נלווה בפרק הזה מערכת קטנה ודמיונית לגמרי: צוות שבונה מסווג (classifier), תוכנה שמשייכת כל קלט לאחת מכמה קטגוריות. המסווג מקבל תמונה של בצק ומכריע: "תפח" או "לא תפח". הצוות והתמונות הם המחשה בלבד.
\section{מאיפה באות הדוגמאות}
הצוות מתחיל באיסוף. אלפי תמונות מצטברות: חלקן צולמו בידי חברי הצוות, חלקן נשלחו מחברים. כבר בשלב הזה נקבע הרבה. לפי מאמר "Datasheets for Datasets", מאפייני מאגר הנתונים משפיעים באופן יסודי על התנהגות המודל, ומודל לא צפוי לתפקד היטב בשטח אם הסביבה שבה הוא פועל אינה תואמת את הנתונים שעליהם אומן ונבדק \cite{S-ch02-04}.
ובכל זאת, כשטימניט גברו ועמיתיה כתבו את המאמר ב-2018, לא היה בקהילת למידת המכונה תהליך סטנדרטי לתיעוד מאגרי נתונים. הם הציעו לצרף לכל מאגר "דף נתונים" (datasheet), בהשראת הדפים שמלווים כל רכיב אלקטרוני \cite{S-ch02-03}. השאלות שבדף נשמעות פשוטות, ודווקא משום כך הן חושפות: מי אסף את הנתונים ובאיזה מנגנון, כמה שולם לו, באיזו שיטה נבחרו הדוגמאות, האם המאגר הוא מדגם, והאם הוא מייצג \cite{S-ch02-04}.
\section{ניקוי: לסלק את מה שמטעה}
אחרי האיסוף מגיע שלב פחות זוהר. ניקוי נתונים, במילים פשוטות, הוא חיפוש אחר מה שעלול להטעות את המכונה, והחלטה מה לעשות בו. קורס המבוא של גוגל ללמידת מכונה מונה כמה חשודים קבועים: ערכים חסרים, דוגמאות כפולות, ערכים שגויים ותוויות שגויות \cite{S-ch02-13}. אצל הצוות זה נראה כך: תמונה שהועלתה פעמיים, תמונה בלי פרטים על מקורה, ועוגה שהתגנבה לתיקייה של לחמים.
אותו מקור מזכיר שאי אפשר לנקות מאגר מכל הרעש \cite{S-ch02-13}. המטרה היא נתונים שאפשר לסמוך עליהם, לא נתונים מושלמים. וכיוון שכל הסרה וכל תיקון משנים את מה שהמודל יראה, דף הנתונים שואל גם אילו פעולות עיבוד בוצעו, והאם הנתונים הגולמיים נשמרו בצד \cite{S-ch02-04}.
\section{תוויות: מי מחליט מה נכון?}
למסווג כזה, תמונה לבדה אינה מלמדת דבר. לכל תמונה צריך להצמיד תווית (label), כלומר את התשובה שהמודל אמור לחזות. וכאן מתברר שהמילה "נכונה" חמקמקה. שני מתייגים מביטים באותו בצק: האחד כותב "תפח", והשנייה מהססת.
אותו קורס, בשיעור על תוויות, ממליץ לא לסמוך בעיניים עצומות על מתייגים אנושיים: הם טועים, לפעמים צריך כמה מהם, וכדאי לתייג בעצמכם כאלף דוגמאות ולהשוות \cite{S-ch02-12}. וכשהמתייגים חלוקים, אין להניח שהתיוג שלכם הוא הנכון, במיוחד בשיפוטים ערכיים \cite{S-ch02-12}.
מלכודת נוספת היא תווית עקיפה (proxy label): במקום למדוד את מה שרוצים, מודדים משהו קרוב. לדוגמה, הצוות עלול להחליט ש"תפח" פירושו "האופה נתן לתוצאה חמישה כוכבים". אבל דירוג מעיד על שביעות רצון, לא בהכרח על תפיחה. תווית עקיפה היא תמיד פשרה, והמודל מועיל רק במידה שהקשר בינה לבין מה שרוצים לחזות חזק \cite{S-ch02-12}.
\begin{behindbox}{שלוש ערימות, אחת נעולה}
אחרי הניקוי והתיוג, הצוות מחלק את התמונות לשלוש ערימות \cite{S-ch02-10}:
קבוצת אימון (training set): ממנה המודל לומד.
קבוצת אימות (validation set): עליה משווים גרסאות ומכווננים היפר-פרמטרים, הגדרות שהאדם קובע ולא המודל.
קבוצת בדיקה (test set): נעולה במגירה, ונפתחת רק לבדיקה הסופית \cite{S-ch02-09}.
קורס המבוא מציג, כדוגמה ולא ככלל, חלוקה של 70\% לאימון ו-15\% לכל אחת מהאחרות \cite{S-ch02-10}. למה לא להסתפק בשתיים? כי כיוונון לפי קבוצת הבדיקה מזרים אל המודל ידע עליה \cite{S-ch02-09}. המחיר: פחות דוגמאות ללמידה, ותוצאה שתלויה במזל של החלוקה. אחת התשובות היא אימות צולב (k-fold): מחלקים את קבוצת האימון ל-k חלקים, בכל סבב לומדים מכולם חוץ מאחד, בוחנים על הנותר, ומדווחים על הממוצע \cite{S-ch02-09}.
\end{behindbox}
\section{הכללה: המבחן שלא ראו לפניו}
למה כל הטקס הזה? כי מודל שנבחן על אותן תמונות שעליהן התאמן יכול להיראות מבריק ולהיות חסר תועלת. התיעוד של הספרייה scikit-learn מתאר זאת בחדות: מודל שפשוט חוזר על התוויות שראה יקבל ציון מושלם, ולא יחזה שום דבר שימושי על נתונים חדשים \cite{S-ch02-09}. לזה קוראים התאמת-יתר (overfitting): שינון של קבוצת האימון, צמוד כל כך שהמודל נכשל מול דוגמאות חדשות \cite{S-ch02-11}. ההפך הוא הכללה (generalization): היכולת לחזות היטב גם על מה שלא נראה מעולם \cite{S-ch02-11}.
הכללה אינה מובטחת. יש תנאים שבהם אפשר לצפות לה: הדוגמאות בלתי תלויות זו בזו ומגיעות מאותה התפלגות, הנתונים אינם משתנים משמעותית לאורך זמן, ושלוש הערימות דומות זו לזו סטטיסטית \cite{S-ch02-11}. ויש גם שחיקה: ככל שבודקים יותר פעמים מול אותה קבוצת בדיקה, המודל מתאים את עצמו אליה, והיא מאבדת מערכה כמבחן \cite{S-ch02-10}. איך מודדים הצלחה בפועל? לשם נגיע בפרק 5.
\begin{figure}[htbp]
\centering
% Figure 2.1: the path of the examples, from collection to the final test (conceptual, no data).
\begin{bookdiagram}[x=1cm,y=1cm]
  \tikzset{st/.style={blk,draw=turquoise,text width=2.75cm,minimum height=11mm}}
  \node[st] (raw)   at (13.6,0)  {דוגמאות גולמיות};
  \node[st] (clean) at (9.6,0)   {ניקוי};
  \node[st] (lab)   at (5.6,0)   {תיוג};
  \node[st] (split) at (1.6,0)   {חלוקה לאימון, אימות ובדיקה};
  \node[st] (train) at (1.6,-2.2) {אימון};
  \node[st] (tune)  at (5.6,-2.2) {כיוונון מול קבוצת האימות};
  \node[st,draw=midnight,line width=1.2pt] (final) at (11.2,-2.2) {בדיקה סופית, חד-פעמית};
  \node[blk,fill=midnight,text=white,text width=3.4cm,minimum height=11mm] (test) at (11.2,-5.0) {קבוצת הבדיקה};
  \draw[flow] (raw) -- (clean);
  \draw[flow] (clean) -- (lab);
  \draw[flow] (lab) -- (split);
  \draw[flow] (split) -- (train);
  \draw[flow] (train) -- (tune);
  \draw[flow] (tune) -- (final);
  \draw[flow] (split.west) -- ++(-0.45,0) |- node[pos=0.75,below,font=\footnotesize] {מופרדת מראש} (test.west);
  % the wall, with the lock where the single legitimate arrow passes
  \draw[midnight,line width=3pt] (8.6,-3.65) -- (14.2,-3.65);
  \draw[midnight,line width=3pt] (8.6,-3.65) -- (8.6,-6.0);
  \draw[flow,draw=midnight] (test.north) -- node[right=4mm,font=\footnotesize] {פעם אחת, בסוף} (final.south);
  \node[fill=gold,draw=midnight,rounded corners=0.6mm,minimum width=5mm,minimum height=4mm,inner sep=0pt] (lock) at (11.2,-3.65) {};
  \draw[midnight,line width=1pt] ([xshift=-1.4mm]lock.north) arc (180:0:1.4mm);
  \draw[loop,line width=1.2pt] (test.west) to[out=180,in=-90] node[pos=0.55,below left,font=\footnotesize\headingfont,text=violet] {זליגה: הדרך האסורה} (tune.south);
  \node[text=violet,font=\headingfont\Large] at (8.6,-4.55) {\(\times\)};
\end{bookdiagram}

\caption{מסלול הדוגמאות, מהאיסוף ועד הבדיקה}\label{fig:2.1}
\figurenote{מקור: מסגרת הספר, על בסיס \cite{S-ch02-08}, \cite{S-ch02-09} ו-\cite{S-ch02-10} (תרשים מושגי, בלי נתונים)}
\end{figure}
\section{זליגה: כשהשאלות דולפות לפני הבחינה}
דמיינו מבחן בגרות שהשאלות שלו דלפו יום קודם. הציונים יזנקו, אבל הם כבר לא אומרים דבר על הידע של התלמידים. בלמידת מכונה קוראים לזה זליגת נתונים (data leakage): שימוש, בזמן בניית המודל, במידע שלא יהיה זמין ברגע החיזוי. התוצאה היא הערכה אופטימית מדי, וביצועים חלשים יותר על נתונים חדשים באמת \cite{S-ch02-08}. שחר קאופמן, סהרון רוסט ועמיתיהם מזכירים שהזליגה נחשבת לאחת מעשר הטעויות המובילות בכריית נתונים, ומציעים למנוע אותה בהפרדה קפדנית בין הנתונים שמהם לומדים לבין אלה שעליהם חוזים \cite{S-ch02-06}.
אצל הצוות הדמיוני, הדליפה מגיעה מהדלת האחורית: עותק כפול אחד חמק מהניקוי ונחת פעם בערימת האימון ופעם בערימת הבדיקה. הקורס של גוגל מזהיר בדיוק מזה: בקבוצת בדיקה טובה אין אף דוגמה משוכפלת מקבוצת האימון, כי המבחן ההוגן היחיד הוא מול דוגמאות חדשות ולא מול העתקים \cite{S-ch02-10}.
הזליגה יכולה להיות עדינה עוד יותר. התיעוד של scikit-learn מציג הדגמה סינתטית, לא מקרה מהעולם: 200 דוגמאות, 10,000 מאפיינים אקראיים ותוויות אקראיות. אין שם שום דבר ללמוד, והדיוק הצפוי הוא של ניחוש, כ-0.5. כשבחרו את המאפיינים "המבטיחים" מתוך כל הנתונים לפני החלוקה, דווח דיוק של 0.76. כשחילקו קודם ובחרו אחר כך, הדיוק חזר ל-0.5 \cite{S-ch02-08}. המודל "גילה" דפוס שלא היה קיים מעולם. הכלל שהתיעוד מנסח: לחלק לאימון ולבדיקה לפני כל עיבוד מקדים, ולעולם לא להתאים (fit) דבר על נתוני הבדיקה \cite{S-ch02-08}.
\section{הטיה: מחברת שנכתבה במטבח אחד}
נחזור למטבח של הקוראת. נניח שכל מה שהיא יודעת על אפייה רשום במחברת אחת, שנכתבה במטבח אחד, מול תנור אחד. בבית שלה המחברת מצוינת; אצל השכנה היא עלולה להטעות. כך גם הצוות: אם כמעט כל התמונות צולמו בשני מטבחים מוארים היטב, המסווג עלול למעוד מול תמונה שצולמה באור עמום של ערב.
לתופעה הזאת קוראים הטיה (bias) בנתונים, ואין לבלבל אותה עם ההטיה שנפגוש בפרק 3, שהיא מספר בתוך הרשת. השיעור על הוגנות באותו קורס מבחין בכמה סוגים, ובהם הטיית בחירה (selection bias), כשהדוגמאות לא נבחרו כך שישקפו את ההתפלגות האמיתית, והטיה היסטורית (historical bias), כשנתוני העבר נושאים אי-שוויון שהיה קיים \cite{S-ch02-14}. מודלים עלולים לשחזר הטיות כאלה ואף להגביר אותן \cite{S-ch02-04}.
המקרה המתועד הבולט, מחקר שנקרא Gender Shades, מגיע מ-2018. ג'וי בואולמוויני וטימניט גברו בחנו שלושה מסווגי מגדר מסחריים, של Microsoft, IBM ו-Face++. קודם הראו ששני מאגרי בנצ'מרק (benchmark, מאגר מבחן מקובל להשוואה) מוכרים, IJB-A ו-Adience, מורכבים ברובם מבהירי עור: 79.6\% ו-86.2\% בהתאמה \cite{S-ch02-01}. לכן בנו מאגר בדיקה משלהן, PPB: 1,270 חברי פרלמנט משלוש מדינות באפריקה ומשלוש באירופה. גוון העור תויג בסולם בן שש דרגות, ורופאת עור מוסמכת קבעה את התוויות הסופיות. המחברות עצמן מציינות שתיוג בינארי של מגדר הוא צמצום שאינו לוכד את המורכבות \cite{S-ch02-01}.
התוצאות על מאגר PPB: אצל נשים כהות עור שיעור השגיאה נע בין 20.8\% ל-34.7\%, לפי המסווג, ואצל גברים בהירי עור בין 0.0\% ל-0.8\% בלבד. נשים כהות עור היו 21.3\% מהמאגר, אבל 61.0\% עד 72.4\% מהשגיאות \cite{S-ch02-01}. הלקח שהספר גוזר מכך: ציון כולל נאה יכול להסתיר כשל עמוק אצל קבוצה אחת, שמתגלה רק כשבודקים כל קבוצה בנפרד.
ומה גרם לפער? כאן נדרשת זהירות. החברות לא פרסמו על אילו נתונים אומנו המסווגים, והמחברות מעלות השערה ולא קביעה: ייתכן שעור כהה מתואם עם מאפייני פנים או עם נורמות של הצגת מגדר שהיו מיוצגים פחות בנתוני האימון \cite{S-ch02-01}.
\begin{mythbox}
\mythlabel אם רק נאסוף מספיק נתונים, ההטיה תתמוסס מעצמה.
\tcblower
\realitylabel הטיית בחירה נובעת מהאופן שבו נבחרו הדוגמאות, לא ממספרן \cite{S-ch02-14}. אם הצוות יכפיל פי עשרה את התמונות משני אותם מטבחים, הוא יקבל מחברת עבה יותר, שנכתבה באותו מטבח. ב-Gender Shades, מה שחשף את הפער היה הרכב ולא גודל: מאגר הבדיקה PPB מאוזן כמעט לגמרי, 53.6\% בהירי עור ו-46.4\% כהי עור \cite{S-ch02-01}. הוא לא תיקן אף מסווג, אבל בעזרתו אפשר היה לראות היכן הם נכשלים. כמות עונה על השאלה "כמה דוגמאות יש?", והטיה שואלת "מי מיוצג, ואיך?". לכן בקבוצת בדיקה טובה גודל לבדו אינו מספיק: היא צריכה גם לייצג את העולם שבו המודל יפעל \cite{S-ch02-10}.
\end{mythbox}
\section{בחזרה לערב ההוא}
הסרט רץ קדימה. הקוראת עדיין מול המסך, והמילים הראשונות של התשובה רק מתחילות להופיע. עכשיו אנחנו יודעים שמאחוריה עומדות החלטות של אנשים שהיא לעולם לא תפגוש: מה לאסוף, מה להשליך, איך לתייג, ומה לנעול במגירה עד הרגע האחרון.
\begin{thinkbox}
דמיינו שלמערכת שהקוראת משוחחת איתה היה מצורף דף נתונים. איזו שאלה הייתם רוצים לקרוא בו ראשונה: מי אסף את הדוגמאות שמהן למדה, מה הוסר מהן בניקוי, או מי החליט מהי תשובה "טובה"? ואם התשובה לשאלה שחשובה לכם ביותר הייתה "לא ידוע", האם הייתם סומכים על העצה באותה מידה?
\end{thinkbox}
נשארת עוד תעלומה. דוגמאות הן תמונות, מילים ומשפטים, אבל המכונה אינה רואה אף אחד מהם. בפרק הבא נרד שכבה נוספת אל מתחת למסך, ונגלה איך הדוגמאות הופכות למספרים, ואיך המספרים מתחילים להכריע.
\begin{insightsbox}
\begin{itemize}
\item כל מערכת לומדת נשענת על שרשרת של החלטות אנושיות: מה נאסף, מה נוקה, ואיך תויג. דף נתונים הוא דרך לחשוף את השרשרת הזאת \cite{S-ch02-03,S-ch02-04}.
\item תווית היא שיפוט, לא עובדה טבעית. מתייגים טועים וחלוקים, ותווית עקיפה מודדת רק משהו קרוב למה שרוצים \cite{S-ch02-12}.
\item קבוצת הבדיקה נשמרת נעולה עד הסוף, כי רק דוגמאות שהמודל לא פגש מעידות על הכללה ולא על שינון \cite{S-ch02-09,S-ch02-11}.
\item זליגה מייצרת הצלחה מדומה: אפילו על נתונים אקראיים לגמרי, עיבוד לפני החלוקה העלה את הדיוק המדווח מרמת ניחוש ל-0.76 \cite{S-ch02-08}.
\item הטיה נובעת מהאופן שבו הנתונים נבחרו וממה שהם משקפים, למשל קבוצה שמיוצגת פחות או אי-שוויון מהעבר, ולא מכמותם; ב-Gender Shades הפער נחשף רק בבדיקה נפרדת של כל תת-קבוצה \cite{S-ch02-01,S-ch02-14}.
\end{itemize}
\end{insightsbox}
